\documentclass[12pt,a4paper]{article}

% ─── Packages ────────────────────────────────────────────────────────────────
\usepackage[top=2.0cm,bottom=2.2cm,left=2.4cm,right=2.4cm]{geometry}
\usepackage[T1]{fontenc}
\usepackage[utf8]{inputenc}
\usepackage{lmodern}
\usepackage{microtype}
\usepackage{xcolor}
\usepackage{titlesec}
\usepackage{graphicx}
\usepackage{enumitem}
\usepackage{hyperref}
\usepackage{fancyhdr}
\usepackage{tcolorbox}
\usepackage{multicol}
\usepackage{tabularx}
\usepackage{colortbl}
\usepackage{hhline}
\usepackage{setspace}
\usepackage{mdframed}
\usepackage{array}
\usepackage{booktabs}

% ─── Rupee Symbol ────────────────────────────────────────────────────────────
\newcommand{\Rup}{\textbf{Rs.\,}}

% ─── Colour Palette ──────────────────────────────────────────────────────────
\definecolor{primaryGreen}{RGB}{22,101,52}
\definecolor{accentGreen}{RGB}{74,222,128}
\definecolor{deepTeal}{RGB}{15,118,110}
\definecolor{softGold}{RGB}{202,138,4}
\definecolor{lightBg}{RGB}{240,253,244}
\definecolor{cardBg}{RGB}{220,252,231}
\definecolor{borderGreen}{RGB}{134,239,172}
\definecolor{textDark}{RGB}{20,83,45}
\definecolor{sectionBar}{RGB}{16,185,129}
\definecolor{orange}{RGB}{194,65,12}
\definecolor{blue}{RGB}{37,99,235}
\definecolor{purple}{RGB}{109,40,217}
\definecolor{rose}{RGB}{190,18,60}
\definecolor{amber}{RGB}{161,98,7}
\definecolor{skyblue}{RGB}{3,105,161}
\definecolor{tierGold}{RGB}{202,138,4}
\definecolor{tierSilver}{RGB}{100,116,139}

% ─── Typography ──────────────────────────────────────────────────────────────
\setlength{\parskip}{5pt}
\setlength{\parindent}{0pt}
\onehalfspacing

% Fixed titleformat: use \rule instead of \titlerule to avoid optional-arg conflict
\titleformat{\section}[block]
  {\large\bfseries\color{primaryGreen}}
  {}{0em}{\colorbox{cardBg}{\parbox{\dimexpr\linewidth-2\fboxsep\relax}{\centering\strut #1\strut}}}[\vspace{2pt}{\color{sectionBar}\rule{\linewidth}{1.2pt}}]

\titleformat{\subsection}[hang]
  {\normalsize\bfseries\color{deepTeal}}
  {\thesubsection.}{0.5em}{}

% ─── Header / Footer ─────────────────────────────────────────────────────────
\pagestyle{fancy}
\fancyhf{}
\fancyhead[L]{\small\color{primaryGreen}\textbf{Kisaan Ki Yash} --- Feature Catalogue}
\fancyhead[R]{\small\color{deepTeal}SIH 2024 Submission}
\fancyfoot[C]{\small\color{gray}\thepage}
\renewcommand{\headrulewidth}{0.4pt}
\renewcommand{\headrule}{\hbox to\headwidth{\color{borderGreen}\leaders\hrule height \headrulewidth\hfill}}

% ─── Custom Boxes ────────────────────────────────────────────────────────────
\tcbuselibrary{skins,breakable,most}

\newtcolorbox{featurebox}[3][]{
  colback=#2!8, colframe=#2, coltitle=white,
  fonttitle=\bfseries\small, title={#3},
  boxrule=1.4pt, arc=6pt, left=8pt, right=8pt, top=6pt, bottom=6pt,
  before upper={\setlength{\parskip}{3pt}}, #1
}

\newtcolorbox{impactbox}{
  colback=softGold!10, colframe=softGold,
  boxrule=1.2pt, arc=5pt, left=8pt, right=8pt, top=6pt, bottom=6pt
}

\newtcolorbox{starbox}{
  colback=tierGold!12, colframe=tierGold,
  boxrule=1.5pt, arc=7pt, left=10pt, right=10pt, top=8pt, bottom=8pt,
  fontupper=\small
}

% ─── Tier Badge macros ───────────────────────────────────────────────────────
\newcommand{\freebadge}{\textcolor{white}{\colorbox{primaryGreen}{\footnotesize\textbf{FREE}}}}
\newcommand{\saathibadge}{\textcolor{white}{\colorbox{blue}{\footnotesize\textbf{\Rup49/mo}}}}
\newcommand{\probadge}{\textcolor{white}{\colorbox{purple}{\footnotesize\textbf{\Rup89/mo}}}}
\newcommand{\bizzbadge}{\textcolor{white}{\colorbox{amber}{\footnotesize\textbf{\Rup499/mo}}}}
\newcommand{\newbadge}{\textcolor{white}{\colorbox{rose}{\footnotesize\textbf{NEW}}}}

% ─── Hyperref ────────────────────────────────────────────────────────────────
\hypersetup{
  colorlinks=true, linkcolor=primaryGreen, urlcolor=deepTeal,
  pdftitle={Kisaan Ki Yash -- Feature Catalogue},
  pdfauthor={Kisaan Ki Yash Team}
}

% ══════════════════════════════════════════════════════════════════════════════
\begin{document}
% ══════════════════════════════════════════════════════════════════════════════

% ─── Cover ───────────────────────────────────────────────────────────────────
\begin{titlepage}
\pagecolor{lightBg}
\centering
\vspace*{1.8cm}

{\color{primaryGreen}\rule{\linewidth}{2.5pt}}
\vspace{0.6cm}

{\Huge\bfseries\color{primaryGreen} Kisaan Ki Yash}\\[0.3cm]
{\large\color{deepTeal}\itshape ``The Glory of the Farmer''}\\[0.5cm]
{\color{sectionBar}\rule{0.4\linewidth}{1pt}}\\[0.8cm]

{\LARGE\bfseries\color{textDark} Complete Feature Catalogue}\\[0.3cm]
{\normalsize\color{gray} From micro-alerts to AI decision intelligence --- every feature, every tier, every impact}

\vspace{1.5cm}

\begin{tcolorbox}[colback=cardBg,colframe=primaryGreen,arc=10pt,boxrule=1.8pt,width=0.82\linewidth,center]
\centering
\textbf{\large\color{primaryGreen}Tier Legend}\\[8pt]
\begin{tabular}{cccc}
\freebadge & \saathibadge & \probadge & \bizzbadge \\[4pt]
Free Forever & Kisan Saathi & Kisan Pro & FPO / Business
\end{tabular}\\[8pt]
\newbadge\ \ {\small = New feature being integrated before final submission}
\end{tcolorbox}

\vspace{1.5cm}

\begin{tabular}{rl}
\textbf{\color{primaryGreen}Platform:} & Kisaan Ki Yash --- 1-km Hyperlocal Agricultural Intelligence \\[4pt]
\textbf{\color{primaryGreen}Coverage:} & Lucknow District, UP (Pilot) $\rightarrow$ Pan-India \\[4pt]
\textbf{\color{primaryGreen}Models:}   & 10 AI/ML Models (M1--M10) \\[4pt]
\textbf{\color{primaryGreen}Live URL:} & \href{https://kisaankiyash.vercel.app}{kisaankiyash.vercel.app} \\
\end{tabular}

\vfill
{\color{primaryGreen}\rule{\linewidth}{2.5pt}}
\end{titlepage}

\pagecolor{white}
\tableofcontents
\newpage

% ══════════════════════════════════════════════════════════════════════════════
\section{MICRO FEATURES --- Always Free}
\textit{Core features accessible to every farmer, zero cost, zero app required}
% ══════════════════════════════════════════════════════════════════════════════

\begin{featurebox}{primaryGreen}{F-01 \quad Hyperlocal 1-km Temperature Forecast \quad \freebadge}
Real-time downscaled temperature for the farmer's exact 1-km grid cell --- not a district average. Powered by Model 1 (Random Forest + AWS calibration). Achieves 39.9\% better accuracy than raw NWP forecasts.

\textbf{Impact:} Farmers know if their field will be 2\textdegree C cooler than the town forecast --- critical for frost-sensitive crops like tomato and mango.
\end{featurebox}

\begin{featurebox}{primaryGreen}{F-02 \quad Daily Rain Yes/No Alert \quad \freebadge}
Simple go/no-go rainfall prediction for the next 24 hours delivered via SMS or WhatsApp. Model 3 hurdle classifier outputs rain probability; a threshold-based binary alert is generated automatically.

\textbf{Impact:} Prevents unnecessary irrigation on a day when rain is coming --- saves water, electricity, and diesel pump costs. One avoided irrigation event saves \Rup80--300 in fuel and labour.
\end{featurebox}

\begin{featurebox}{primaryGreen}{F-03 \quad 3-Day Basic Weather Forecast \quad \freebadge}
Temperature high/low, humidity, and wind speed for 3 days at 1-km resolution, displayed in Hindi on a simple card view. No internet required beyond the initial SMS notification.

\textbf{Impact:} Helps farmers plan field operations (spraying, harvesting, sowing) 3 days in advance, reducing reactive scrambling and labour inefficiency.
\end{featurebox}

\begin{featurebox}{primaryGreen}{F-04 \quad WhatsApp \& SMS Advisory Delivery \quad \freebadge}
Zero-app-install distribution. All advisories are delivered as plain Hindi text via WhatsApp Business API or SMS gateway --- works on any \Rup500 feature phone.

\textbf{Impact:} Eliminates smartphone adoption as a barrier to access. Reaches the bottom 60\% of farmers who do not own smartphones.
\end{featurebox}

\begin{featurebox}{primaryGreen}{F-05 \quad Panchayat-Level Location Tagging \quad \freebadge}
All forecasts are tagged to the farmer's specific panchayat (5 pilot panchayats in Lucknow, expanding to 6,000+ across UP). No GPS required --- farmer registers their village name once during onboarding.

\textbf{Impact:} Replaces generic district forecasts with village-level accuracy. A farmer in Malihabad gets a different forecast from one in Mohanlalganj --- even though both are in Lucknow district.
\end{featurebox}

% ══════════════════════════════════════════════════════════════════════════════
\section{OPERATIONAL FEATURES --- Kisan Saathi (\Rup49/month)}
\textit{Input optimisation and daily farm management intelligence}
% ══════════════════════════════════════════════════════════════════════════════

\begin{featurebox}{blue}{F-06 \quad Soil Moisture Status \quad \saathibadge}
Surface and root-zone soil moisture estimate for the farmer's 1-km grid cell, fusing Sentinel-1 C-band SAR backscatter with ISRIC SoilGrids 250m soil texture data (Model 4 --- M4).

Example output: \textit{``Your field's root zone has 24\% volumetric moisture --- above the irrigation threshold for paddy. Skip irrigation today. Next recommended: Oct 4.''}

\textbf{Impact:} Saves 2--4 unnecessary irrigation events per crop cycle. At \Rup200/event (pump diesel + labour), this is \Rup400--800/season just on avoided irrigation.
\end{featurebox}

\begin{featurebox}{blue}{F-07 \quad Irrigation Scheduling Engine \quad \saathibadge}
Daily irrigation requirement in mm computed using FAO-56 Penman-Monteith evapotranspiration, calibrated with 1-km temperature, humidity, and wind from M1/M2, and soil moisture deficit from M4. Output: \textit{``Irrigate 35 mm tomorrow morning. Next required: Oct 7.''}

\textbf{Impact:} 15--25\% water saving per season. In Uttar Pradesh, where Ganga-plain groundwater is depleting at 0.5 m/year, precision irrigation is both an economic and environmental imperative.
\end{featurebox}

\begin{featurebox}{blue}{F-08 \quad 7-Day Extended Forecast \quad \saathibadge}
Full 7-day hyperlocal forecast at 1-km resolution --- temperature, humidity, rainfall probability, wind speed, and solar radiation. Enables complete weekly farm planning in a single Sunday-morning advisory.

\textbf{Impact:} Farmers can plan labour booking, equipment hire, market trips, and crop operations a full week in advance, reducing last-minute costs by an estimated \Rup500--1,500/week.
\end{featurebox}

\begin{featurebox}{rose}{F-09 \quad \textbf{Weather-Triggered Fertiliser Alert} \quad \saathibadge\quad\newbadge}
The platform continuously monitors the 7-day precipitation forecast for each panchayat. When an optimal fertiliser application window is detected --- soil is moist enough for uptake but rainfall is 4--7 days away (preventing runoff) --- an automatic SMS/push notification is sent.

\vspace{4pt}
\textbf{Example advisory message:}
\begin{mdframed}[backgroundcolor=cardBg,linecolor=primaryGreen,linewidth=1pt,roundcorner=4pt]
\textit{``Alert for Malihabad: Rain is expected in your region in 6 days. This is the optimal window to apply your base dose of DAP/Urea fertiliser. Apply before October 3rd for best root uptake. Avoid application within 48 hours of expected heavy rainfall --- nutrient runoff will waste your entire investment.''}
\end{mdframed}

\vspace{4pt}
The system also issues a \textbf{PROHIBIT alert} when rain is forecast within 48 hours: \textit{``Do NOT apply fertiliser today. Heavy rain forecast within 2 days. Your nutrients will wash away.''}

\textbf{Impact:} Prevents \Rup2,000--8,000/acre loss per season from nutrient runoff caused by fertiliser applied immediately before heavy rain. Also prevents the reverse --- applying in too-dry soil where uptake efficiency drops to below 30\%.

\textbf{Monetisation opportunity:} Agri-input companies (IFFCO, Coromandel, UPL) can sponsor precision-targeted fertiliser alerts as agronomically justified advertising --- revenue neutral or even revenue-positive for the platform.
\end{featurebox}

\begin{featurebox}{blue}{F-10 \quad Pesticide / Spray Timing Advisor \quad \saathibadge}
Validates forecast conditions before recommending pesticide or fungicide spray. Spraying within 4 hours of rain = 100\% chemical washoff = completely wasted input. System confirms a dry window of 6+ hours with wind speed $<$ 15 km/h before issuing a ``spray safe'' advisory.

\textbf{Impact:} Saves \Rup500--1,500/acre per spray event. Across 50,000 farms with 3 spray events per season, this prevents \Rup7.5 crore in annual wasted chemical input --- and reduces soil and water chemical load proportionally.
\end{featurebox}

\begin{featurebox}{blue}{F-11 \quad Crop Stage Tracking (GDD Calendar) \quad \saathibadge}
NDVI and EVI time series from Sentinel-2 tracks crop growth stage (emergence, tillering, booting, heading, maturity) and maps to Growing Degree Day (GDD) accumulation. Farmer receives stage-specific nutrient and protection recommendations.

Example: \textit{``Your wheat is at tillering stage. 47 GDD accumulated. Apply split-dose top-dress nitrogen now for maximum tiller activation.''}

\textbf{Impact:} Prevents split-dose misapplication --- one of the top causes of sub-optimal yield in wheat and paddy. Estimated 8--12\% yield improvement from correctly timed nitrogen application.
\end{featurebox}

% ══════════════════════════════════════════════════════════════════════════════
\section{STRATEGIC FEATURES --- Kisan Pro (\Rup89/month)}
\textit{Season-level protection and market intelligence --- protect the entire crop, not just tomorrow}
% ══════════════════════════════════════════════════════════════════════════════

\begin{featurebox}{purple}{F-12 \quad Extreme Weather Hazard Alerts (48--72h Advance) \quad \probadge}
Model 9 (M9) detects anomalous atmospheric events against 30-year ERA5 climatological baselines and issues location-specific alerts 48--72 hours in advance:
\begin{itemize}[leftmargin=*,noitemsep]
  \item \textbf{Heatwave:} Temperature $>$ 95th percentile for 3+ consecutive days
  \item \textbf{Cold Wave / Frost Risk:} Minimum temperature $<$ 4\textdegree C during Rabi season (wheat, mustard, potato)
  \item \textbf{Prolonged Dry Spell:} Zero rainfall for $>$ 7 consecutive days during Kharif
  \item \textbf{Waterlogging Risk:} Cumulative rainfall $>$ field capacity on detected impermeable-layer soils
  \item \textbf{Loo (hot dry wind):} Extreme VPD index during May--June for horticulture protection
\end{itemize}

\textbf{Impact:} A 48-hour frost alert allows a mango farmer in Malihabad to apply protective smoke, cover young shoots, or arrange micro-irrigation for heat -- preventing 60--100\% crop loss that would otherwise occur overnight. This single alert can save \Rup40,000--1,20,000 per acre for an established orchard.
\end{featurebox}

\begin{featurebox}{purple}{F-13 \quad Crop Yield Forecast \quad \probadge}
End-of-season yield estimate (quintals/hectare) computed using cumulative weather anomalies, vegetation stress index, and irrigation history --- calibrated against district-level Crop Cutting Experiment (CCE) data (Model 7, M7).

\textbf{Impact:} Farmers can forward-sell at better prices when yield outlook is strong. Farmer Producer Organisations can negotiate bulk procurement contracts. Insurance companies receive objective, satellite-verified pre-harvest yield estimates --- eliminating disputes.
\end{featurebox}

\begin{featurebox}{purple}{F-14 \quad Flood \& Waterlogging Risk Map \quad \probadge}
Topographic Wetness Index (TWI) and DEM depression analysis identifies 1-km cells structurally prone to waterlogging. M3 precipitation forecast drives a simple water-balance to issue 24--72 hour inundation risk maps per panchayat (Model 8, M8).

\textbf{Impact:} Paddy farmers in low-lying fields can pre-open drainage channels before flood arrival. Saves 15--40\% yield in high-risk years, particularly in UP's eastern districts (Gorakhpur, Ballia belt).
\end{featurebox}

\begin{featurebox}{rose}{F-15 \quad Hindi Voice Advisory --- IVR Phone Call \quad \probadge\quad\newbadge}
For farmers who cannot read SMS, an automated outbound IVR (Interactive Voice Response) call delivers the same advisory in spoken, colloquial Hindi. Triggered automatically when a critical or extreme-weather alert is issued.

\vspace{4pt}
\textbf{Example IVR script (transliterated):}\\
\textit{``Namaste! Yeh Kisaan Ki Yash ki taraf se ek zaroori sandesh hai. Kal aapke kshetra mein tez barish hone ki sambhavana hai. Kripya aaj hi apni paki fasal ki kutti kaaten aur khule mein rakha anaaj dhak len. Dhanyavaad.''}

\textit{(English: Namaste! This is an important message from Kisaan Ki Yash. Heavy rain is forecast in your area tomorrow. Please harvest your mature crop today and cover any grain kept in the open. Thank you.)}

\textbf{Impact:} Extends the platform's protective intelligence to illiterate and elderly farmers --- the demographic most financially vulnerable to crop loss and least served by existing digital platforms.
\end{featurebox}

\begin{featurebox}{purple}{F-16 \quad PMFBY Insurance Report Auto-Generation \quad \probadge}
Verified M1/M9 weather event records (hailstorm timestamps, flood extent, temperature anomaly duration) are automatically formatted into structured crop damage evidence reports compatible with PMFBY (Pradhan Mantri Fasal Bima Yojana) claim submission requirements.

\textbf{Impact:} Eliminates contested, subjective field inspections. Objective, timestamped, satellite-verified evidence reduces claim settlement time from months to days. Enables farmers to receive insurance compensation before the next sowing season --- when they need it most.
\end{featurebox}

\begin{featurebox}{purple}{F-17 \quad Conformal Prediction Intervals \quad \probadge}
Every temperature and rainfall forecast includes scientifically computed uncertainty bounds at 90\% nominal coverage: \textit{``Temperature: 28.4\textdegree C [27.5\textdegree C -- 29.3\textdegree C]''}. Bounds are honestly withheld (displayed as ``Calibration Limited'') when training data is insufficient --- the system abstains rather than fabricates confidence.

\textbf{Impact:} Farmers and agronomists can see exactly how confident the system is. High uncertainty = defer the decision. Low uncertainty = act immediately. No other agricultural weather platform in India provides this level of calibrated, honest transparency.
\end{featurebox}

% ══════════════════════════════════════════════════════════════════════════════
\section{ENTERPRISE FEATURES --- FPO / Business (\Rup499/month)}
\textit{Portfolio-level intelligence for Farmer Producer Organisations, cooperatives, and agri-businesses}
% ══════════════════════════════════════════════════════════════════════════════

\begin{featurebox}{amber}{F-18 \quad Multi-Farm FPO Dashboard \quad \bizzbadge}
Aggregate view of all enrolled farms in the FPO: panchayat-level heatmaps of soil moisture, yield forecast, and risk alerts. Colour-coded block-level risk summary with one-click drill-down. Export to Excel or PDF for management meetings and board reporting.

\textbf{Impact:} An FPO managing 500 farmer members can identify at a glance which blocks need emergency input intervention vs.\ which are performing on track --- enabling efficient allocation of extension officer resources.
\end{featurebox}

\begin{featurebox}{amber}{F-19 \quad REST API Access \quad \bizzbadge}
Full programmatic access to all 10 model outputs via documented REST API endpoints (FastAPI + OpenAPI spec). All responses include provenance records, model version, and data lineage. Enables embedding into ERP systems, rural bank credit assessment tools, and custom dashboards.

\textbf{Impact:} Agri-input companies, seed companies, rural banks, and state government agencies can embed live farm intelligence into their own digital workflows without rebuilding the data infrastructure.
\end{featurebox}

\begin{featurebox}{rose}{F-20 \quad e-NAM Price Signal Integration \quad \bizzbadge\quad\newbadge}
Crop yield forecast (M7) is cross-referenced with real-time e-NAM (National Agricultural Market) commodity prices to generate optimal harvest-timing recommendations:

\textit{``Current paddy MSP at Lucknow mandi is \Rup2,183/quintal --- 12\% above 3-month average. Your field is 10--14 days from harvest. Consider early harvest to capture this price window before new-season arrivals suppress the market.''}

\textbf{Impact:} Farmers sell at peak prices rather than distress-selling during harvest glut. Estimated 10--25\% income uplift on market-connected farms. Reduces post-harvest losses from delayed selling by enabling proactive logistics planning.
\end{featurebox}

\begin{featurebox}{rose}{F-21 \quad Carbon Credit Verification Reports \quad \bizzbadge\quad\newbadge}
Precision irrigation scheduling (M6) combined with satellite-verified soil moisture records (M4) enables measurement of water savings and reduced groundwater extraction relative to regional baseline. Formatted as Verra / Gold Standard Verified Carbon Standard (VCS) methodology-compatible MRV (Monitoring, Reporting, Verification) documents.

\textbf{Impact:} An FPO managing 1,000 acres of irrigated farmland can potentially generate 200--400 tCO\textsubscript{2}e of verified emissions reductions per year, earning \Rup5--15 lakh/year in voluntary carbon market credits --- a significant additional income stream on top of crop revenue.
\end{featurebox}

% ══════════════════════════════════════════════════════════════════════════════
\section{MASTER DECISION ENGINE --- M10 (All Tiers)}
\textit{The brain of the platform --- synthesises all upstream models into one actionable advisory per panchayat}
% ══════════════════════════════════════════════════════════════════════════════

\begin{tcolorbox}[colback=lightBg,colframe=primaryGreen,arc=8pt,boxrule=1.8pt,
  title={\bfseries\color{white}\large F-22 \quad Master Agricultural Decision Advisory (M10)},
  coltitle=white]

The Agricultural Decision Intelligence Engine (Model 10) aggregates outputs from M1 through M9 into a single, prioritised advisory per panchayat, delivered in plain Hindi. It uses a rule engine with expected-loss minimisation to determine the urgency and type of each recommendation.

\vspace{8pt}

\begin{center}
\small
\renewcommand{\arraystretch}{1.4}
\begin{tabular}{|p{2.8cm}|p{3cm}|p{5cm}|p{2cm}|}
\hline
\rowcolor{primaryGreen}
\textbf{\color{white}Farm Action} & \textbf{\color{white}Status} & \textbf{\color{white}Trigger Condition} & \textbf{\color{white}Urgency} \\
\hline
Irrigate & RECOMMENDED & Soil moisture $<$ threshold, no rain in 5 days & LOW \\
\hline
Irrigate & PROHIBITED & Rain forecast within 24 h & NONE \\
\hline
Spray Pesticide & CONDITIONAL & Dry window $>$ 6 h, wind $<$ 15 km/h & MODERATE \\
\hline
Apply Fertiliser & RECOMMENDED & Rain in 5--7 days, optimal uptake window & HIGH \\
\hline
Apply Fertiliser & PROHIBITED & Rain $<$ 48 h away & HIGH \\
\hline
Sow Crop & CONDITIONAL & Soil temperature $>$ threshold, moisture OK & MODERATE \\
\hline
Harvest & CRITICAL & Rain forecast imminent, delayed harvest = total loss & CRITICAL \\
\hline
Open Field Drains & CRITICAL & Waterlogging risk $>$ 70\% within 24 h & CRITICAL \\
\hline
Frost Protection & CRITICAL & Min temperature $<$ 4\textdegree C within 24 h during sensitive growth stage & CRITICAL \\
\hline
\end{tabular}
\end{center}

\vspace{8pt}
\textbf{Impact:} Replaces the need for a personal agronomist. Every farmer receives the equivalent of a trained crop scientist making real-time, evidence-based decisions for their specific field --- at the cost of less than one cup of tea per day.
\end{tcolorbox}

% ══════════════════════════════════════════════════════════════════════════════
\section{IMPACT SUMMARY}
% ══════════════════════════════════════════════════════════════════════════════

\begin{starbox}
\begin{center}
{\large\bfseries\color{primaryGreen} Projected Impact at Scale --- 50,000 Farmers, Year 1}
\end{center}
\vspace{6pt}
\begin{multicols}{2}
\textbf{Economic Impact:}
\begin{itemize}[leftmargin=*,noitemsep]
  \item \Rup4,000--12,000 saved per farmer per year
  \item \Rup500--1,500 saved per avoided spray event
  \item \Rup2,000--8,000/acre saved from timed fertiliser application
  \item 8--15\% crop loss prevented via early hazard alerts
  \item 10--25\% income uplift from optimal market timing (e-NAM)
  \item \Rup5--15 lakh/year potential carbon credit income per FPO (1,000 acres)
\end{itemize}

\columnbreak

\textbf{Environmental Impact:}
\begin{itemize}[leftmargin=*,noitemsep]
  \item 15--25\% reduction in irrigation water usage per farm
  \item 20\% reduction in fertiliser nutrient runoff per event
  \item 30\% reduction in over-irrigation and waterlogging events
  \item Measurable N\textsubscript{2}O emission reduction from precision N timing
  \item Slowing groundwater depletion in UP aquifer zone
  \item Satellite-verified PMFBY claims reducing crop insurance fraud
\end{itemize}
\end{multicols}
\end{starbox}

\vspace{8pt}

\begin{impactbox}
\textbf{Fertiliser Alert Feature Alone --- Back-of-Envelope Impact:}\\[4pt]
UP produces approximately 60 million metric tonnes of food grain annually. Industry studies suggest 10--15\% of applied fertiliser is lost to runoff from mistimed application (applying immediately before heavy rain). At \Rup25,000/MT average fertiliser cost, even a conservative 5\% improvement in application efficiency across 50,000 farmers with 2 acres each = \textbf{\Rup15 crore in annual input savings} --- and an equivalent reduction in river nitrogen load and nitrous oxide emissions.
\end{impactbox}

% ══════════════════════════════════════════════════════════════════════════════
\section{FEATURE SUMMARY TABLE}
% ══════════════════════════════════════════════════════════════════════════════

\begin{center}
\small
\renewcommand{\arraystretch}{1.35}
\begin{tabular}{|c|p{6.5cm}|c|c|}
\hline
\rowcolor{primaryGreen}
\textbf{\color{white}ID} & \textbf{\color{white}Feature Name} & \textbf{\color{white}Tier} & \textbf{\color{white}New?} \\
\hline
F-01 & Hyperlocal 1-km Temperature Forecast & Free & --- \\
F-02 & Daily Rain Yes/No Alert & Free & --- \\
F-03 & 3-Day Basic Weather Forecast & Free & --- \\
F-04 & WhatsApp \& SMS Advisory Delivery & Free & --- \\
F-05 & Panchayat-Level Location Tagging & Free & --- \\
\hline
F-06 & Soil Moisture Status & \Rup49 & --- \\
F-07 & Irrigation Scheduling Engine & \Rup49 & --- \\
F-08 & 7-Day Extended Forecast & \Rup49 & --- \\
\rowcolor{rose!20}
F-09 & \textbf{Weather-Triggered Fertiliser Alert} & \Rup49 & \textbf{NEW} \\
F-10 & Pesticide / Spray Timing Advisor & \Rup49 & --- \\
F-11 & Crop Stage Growth Calendar (GDD) & \Rup49 & --- \\
\hline
F-12 & Extreme Weather Hazard Alerts (48--72h advance) & \Rup89 & --- \\
F-13 & Crop Yield Forecast (end-of-season) & \Rup89 & --- \\
F-14 & Flood \& Waterlogging Risk Map & \Rup89 & --- \\
\rowcolor{rose!20}
F-15 & \textbf{Hindi Voice Advisory via IVR} & \Rup89 & \textbf{NEW} \\
F-16 & PMFBY Insurance Report Auto-Generation & \Rup89 & --- \\
F-17 & Conformal Prediction Intervals (honest uncertainty) & \Rup89 & --- \\
\hline
F-18 & Multi-Farm FPO Portfolio Dashboard & \Rup499 & --- \\
F-19 & REST API Programmatic Access & \Rup499 & --- \\
\rowcolor{rose!20}
F-20 & \textbf{e-NAM Market Price Signal Integration} & \Rup499 & \textbf{NEW} \\
\rowcolor{rose!20}
F-21 & \textbf{Carbon Credit MRV Verification Reports} & \Rup499 & \textbf{NEW} \\
\hline
\rowcolor{cardBg}
F-22 & Master Agricultural Decision Advisory (M10) & All Tiers & --- \\
\hline
\end{tabular}
\end{center}

\vspace{0.6cm}

\begin{center}
{\color{sectionBar}\rule{0.6\linewidth}{1.5pt}}\\[10pt]
{\large\bfseries\color{primaryGreen} Kisaan Ki Yash --- When science reaches the field,}\\
{\large\bfseries\color{primaryGreen} that is the farmer's glory.}\\[8pt]
{\normalsize\color{deepTeal}\href{https://kisaankiyash.vercel.app}{kisaankiyash.vercel.app}}
\end{center}

\end{document}
